\section{Related Work}

Our work sits at the intersection of three research threads: empirical studies of benchmark overfitting and saturation, qualitative critiques of ML evaluation practices, and the application of mathematical growth models to characterize technological progress.

\subsection{Benchmark Overfitting and Generalization Gaps}

\citet{recht2019imagenet} demonstrated an 11--14 percentage point accuracy gap between ImageNet's original test set and a new independent collection (ImageNet-v2), establishing empirically that models overfit to benchmark-specific features.
This cross-sectional approach is rigorous but fundamentally limited for our purposes: it requires collecting a new independent test set for each benchmark, making it resource-intensive and impossible to apply retrospectively at scale.
\citet{engstrom2020identifying} and \citet{gururangan2018annotation} examined dataset artifacts as alternative signals, focusing on static structural properties rather than temporal progression.
\citet{ethayarajh2020utility} proposed a utility measure for benchmarks but did not address temporal saturation.
In contrast, our approach requires only historical leaderboard timeseries already available publicly---no new data collection required.

\subsection{Qualitative Saturation Discourse and Goodhart's Law}

\citet{wang2019superglue} explicitly motivated SuperGLUE by GLUE's saturation; \citet{srivastava2022bigbench} similarly motivated BIG-bench by SuperGLUE's saturation.
Both treat the saturation event as a community judgment rather than providing a formal detection criterion.
\citet{goodhart1975problems}'s observation---``when a measure becomes a target, it ceases to be a good measure''---provides the theoretical framing for why benchmarks saturate.
\citet{manheim2019categorizing} operationalizes this in the ML context.
However, neither work provides a quantitative temporal characterization.
\citet{bender2021stochastic} and \citet{kapoor2023leakage} offer qualitative critiques without automated detection machinery.
The gap is a formal, reproducible measurement procedure---which we address.

\subsection{Growth Models for Technological Progress}

Logistic growth models have a long history in modeling processes with initial rapid growth followed by a carrying-capacity asymptote: population dynamics \citep{verhulst1838notice}, technology adoption \citep{bass1969new_product}.
In the ML context, \citet{kaplan2020scaling} studied power-law scaling trends but focused on aggregate trends rather than per-benchmark temporal saturation.
We are not aware of prior work applying logistic growth models to benchmark leaderboard timeseries for saturation detection.
The AIC framework \citep{burnham2002model} is well-established in ecology; we transfer it to benchmark dynamics.

\textbf{Positioning}: Our work is longitudinal (temporal) rather than cross-sectional; automated and benchmark-agnostic rather than requiring domain-specific new test sets; and formally information-theoretic rather than qualitative.
